% Source of truth for factual CV content.

% MLOPS
% \section{Professional Summary}
% End-to-end Machine Learning Engineer with solid MLOps experience and a strong research foundation. Proven track record in deploying production-scale ML systems, improving accuracy, reducing latency, and driving real-time model performance. Adept at collaborating across science, engineering and product teams to build impactful solutions.


\section{Professional Summary}
Machine Learning \& Computer Vision Engineer with research experience since 2020 and 3+ years delivering production systems for medical imaging. Core stack: \textbf{Python} and \textbf{PyTorch}; experience spans 2D--3D registration, X-ray image analysis, and full-cycle ML model development and deployment.

\section{Experience}
\resumeSubHeadingList
  \resumeSubheading
      {Machine Learning \& Computer Vision Engineer}{Jun 2023 -- Jul 2026}
      % Public company context, not a claim about the candidate's own dataset or clients.
      % Source: https://www.raylytic.com/en/references/patient-registry-software/
      {RAYLYTIC GmbH -- Medical imaging \& clinical-trial software; \projectlink{https://www.raylytic.com/en/references/patient-registry-software/}{2,000+ patients across 50+ sites}}{Full-time}
      \resumeSubHeadingList
          % Candidate confirmed on 2026-10-29: typical clinical-study volume was approximately 2,000--3,000 images per study.
          \resumeItem{\textbullet\ Built an end-to-end \textbf{2D--3D registration} pipeline for implant pose estimation from X-rays, including mesh processing and alignment, used in clinical studies with typically \textbf{2,000--3,000 images per study}; improved acceptance rate by 45\% and reduced error by 15\%; deployed with GPU-accelerated \textbf{Docker} and \textbf{Prefect}}
          \resumeItem{\textbullet\ Implemented an \textbf{MLOps} pipeline with data versioning via \textbf{DVC}, \textbf{MLflow} experiment tracking, and automated evaluation}
          \resumeItem{\textbullet\ Delivered production \textbf{segmentation, detection, and classification} models for X-ray imaging; improved evaluation metrics (mAP, Dice, ROC) by 20--50\% over baseline and reduced training time 8-fold}
          \resumeItem{\textbullet\ Trained a DINOv3 image encoder self-supervised on X-ray data, then used it in a customized ViTDet architecture for supervised instance detection training on labeled data}
          \resumeItem{\textbullet\ Migrated segmentation model serving to a centralized \textbf{Triton Inference Server} REST API, cutting latency by 50\%}
          
          \resumeItem{\textbullet\ Developed 3D instance segmentation and registration for light-sheet microscopy images in a \projectlink{https://www.photonikforschung.de/projekte/lebenswissenschaften/projekt/photiomics.html}{liver cancer research project}.}
      \resumeSubHeadingListEnd
  
  \resumeSubheading
      {Freelance Computer Vision Engineer}{Feb 2026 -- Mar 2026}
      {Smoothr.de / Sterlix GmbH}{Contract}
      \resumeSubHeadingList
          \resumeItem{\textbullet\ Developed depth-map--based camera obstruction detection for real-time retail anti-theft systems}
      \resumeSubHeadingListEnd
      
  \resumeSubheading
      {Research Engineer}{Jan 2022 -- Mar 2023}
      {Data Science Research Group}{Part-time}
      \resumeSubHeadingList
          % Confirmed context: learned distributed training on the university HPC cluster and introduced PyTorch Lightning to the research group.
          \resumeItem{\textbullet\ Developed a demo ML model for bio-signal time-series classification using a high-performance computing cluster for pain-level assessment in surgical scenarios with non-responsive patients}
          \resumeItem{\textbullet\ Designed data pipelines and training workflows for a ticket-routing classifier, transforming unstructured text into structured features and streamlining dataset updates and model evaluation for \projectlink{https://infai.org/projekt-kmi/}{industry partners}.}
      \resumeSubHeadingListEnd

  \vspace{2pt}
  \textbf{Internships}
  \resumeSubHeadingList
      \resumeItem{\textbullet\ {ML Thesis Student at RAYLYTIC GmbH} -- (Oct 2022 -- Feb 2023): Developed anatomical landmark localization model}
      \resumeItem{\textbullet\ {ML Intern at Prudsys GmbH} -- (Dec 2021 -- Feb 2022): Automated analytic workflows and data preprocessing pipelines}
      \resumeItem{\textbullet\ {Autonomous Driving Intern at IAV GmbH} -- (Sep 2021 -- Nov 2021): Investigated inverse {reinforcement learning} for self-driving}
      \resumeItem{\textbullet\ {Student Research Assistant at WHZ} -- (Jan 2020 -- Dec 2021): Developed a \projectlink{https://www.youtube.com/watch?v=Xa7rivGy6WQ}{semantic segmentation model for automated tissue-density analysis of 3D lung CT scans}; developed an {autoML} library and predicted patient admission to the ICU during COVID.}
  \resumeSubHeadingListEnd
\resumeSubHeadingListEnd

\section{Technical Skills}
\resumeSubHeadingList
  \resumeItem{\textbf{Programming Languages}: \textbf{Python}, C++17, TypeScript, Java, MATLAB}
  \resumeItem{\textbf{ML \& CV}: \textbf{PyTorch}, PyTorch Lightning, MONAI, TensorFlow, JAX, Detectron2, MMEngine/MMCV/MMDetection, Albumentations, scikit-learn, XGBoost, OpenCV, ITK, pyclesperanto, NumPy, pandas, CVAT, Hugging Face Transformers, timm}
  \resumeItem{\textbf{Experimentation \& Tracking}: Jupyter Notebook/JupyterLab, MLflow, TensorBoard, Optuna, DVC, Matplotlib/seaborn/Plotly}
  % Candidate confirmed on 2026-11-01: uses conda and mamba; used AWS for image classification model hosting.
  \resumeItem{\textbf{Infrastructure \& Deployment}: Linux, Docker, Kubernetes, conda, mamba, Prefect, CI/CD (GitLab, Bitbucket), AWS, Triton Inference Server, TensorRT}
  % Confirmed by candidate: Python/FastAPI segmentation and detection services, with unit and performance tests in CI; uses Codex and Claude Code.
  \resumeItem{\textbf{Backend \& Dev Tools}: Flask, REST, FastAPI, OpenAPI, FHIR, Git, SQL, MQTT, Codex, Claude Code}
\resumeSubHeadingListEnd

\section{Education}
\resumeSubHeadingList
  \resumeSubheading
      {B.Sc.\ Data Science with distinction (1.1 / 1.0)}{Oct 2019 -- Apr 2023}
      {University of Applied Sciences Zwickau}{Zwickau, Germany}
      % \textbullet\ \small{Thesis: Anatomical landmark localization for lumbar segmental angle estimation}
\vspace{-2pt}
  \resumeSubheading
      {B.Sc.\ and M.Sc.\ Civil Engineering (M.Sc.\ wGPA: 90.12/100)}{Sep 2014 -- Jun 2020}
      {Odesa State Academy of Civil Engineering and Architecture}{Odesa, Ukraine}
      % \textbullet\ \small{Thesis: Dynamics and stability of rectangular orthotropic plates}
\resumeSubHeadingListEnd

\section{Publications \& Awards}
\resumeSubHeadingList
  \resumeItem{\textbf{EFORT 2026}: Validation of a 2D-3D Image Registration Method for TKA Implant Pose Estimation. Pysanyi et al.}
  \resumeItem{\textbf{DWG 2023}: Automatic determination of lumbar segmental lordosis based on an AI algorithm. Pysanyi et al.}
  \resumeItem{\textbf{1st Place}: National Strength of Materials Olympiad (Ukraine, 2018)}
\resumeSubHeadingListEnd

\section{Languages}
\resumeSubHeadingList
  \resumeItem{German (C2; TestDaF C1, 2019), English (C2, EF SET Certificate)}
\resumeSubHeadingListEnd
